\section{结论}

本文建立了适用于本题长圆柱药材的径向传热传质模型，并按题目要求逐步处理固定尺寸、变物性、全域达标和径向收缩四种情形。

问题一中，预热阶段采用温度—水分解耦模型。$1800\,\mathrm{s}$ 时中心温度为 $33.5753\,{}^\circ\mathrm{C}$、表面温度为 $36.7856\,{}^\circ\mathrm{C}$；中心水分浓度为 $2.5500\,\mathrm{kg/kg}$、表面为 $1.5102\,\mathrm{kg/kg}$。结果表明热量已向内部传播，而水分变化仍主要集中在表面附近。

问题二在固定半径上引入附录3的变物性和温度—水分参数耦合，给出了前三小时指定时刻和位置的温度、水分场。数值结果显示温度响应快于水分响应，轴线水分浓度在早期明显滞后；网格、容差和积分平衡检验支持表格结果的稳定性。

问题三以全域最大水分浓度为判停量，在固定半径和附录3物性下得到临界时刻约 $57.17\,\mathrm{h}$。终点最大值位于轴线，表面水分浓度约为 $0.0525\,\mathrm{kg/kg}$。该结果说明表面先达到阈值不能替代全域判定。

问题四将附件2半径轨迹和附录4物性纳入材料坐标模型，得到全域最大水分浓度首次降至 $0.15\,\mathrm{kg/kg}$ 的临界时刻为 $182956.15\,\mathrm{s}$，即约 $50.82\,\mathrm{h}$；终点半径为 $1.2000\,\mathrm{cm}$，表面水分浓度约为 $0.0525\,\mathrm{kg/kg}$。在给定模型条件下，收缩缩短传递距离，但附录4物性会改变整体传质能力，二者共同决定最终时间。

网格收敛、时间积分稳定性、积分守恒、零扩散极限、固定域退化复现和代表时刻边界根扫描共同支持本文数值实现的可靠性与模型内部自洽性。由于模型仍依赖径向一维、长度不变、环境末值延拓及未显式计入潜热等假设，后续实际工艺应用应结合端面数据和独立实验进行校准。
